
\documentclass{beamer}

% Tema do slide
\usetheme{Madrid}

% Pacotes
\usepackage[utf8]{inputenc}
\usepackage[T1]{fontenc}
\usepackage[brazilian]{babel}
\usepackage{amsmath, amssymb, amsthm}
\usepackage{tikz}
\usetikzlibrary{automata, positioning, arrows.meta}
\usepackage{booktabs}

% Comandos personalizados
\newcommand{\Lt}{\mathsf{L}}
\newcommand{\Nm}{\mathsf{N}}
\newcommand{\Hf}{\mathsf{H}}
\newcommand{\LAJ}{\Lt_{\mathrm{AJ}}}
\newcommand{\lang}[1]{\mathcal{L}(\mathrm{#1})}

% Informações da apresentação

\title{Validador de Placas Veiculares Brasileiras}
\subtitle{Via Autômato Finito Determinístico}
\author{
Celso Vinicius Souza Silva \and
Daniel Paulino da Silva Neto \and
João Gabriel da Silva Nemezio \and
Wlademir Pontes Xavier
}
\date{\today}

\begin{document}

% CAPA

\begin{frame}
    \titlepage
\end{frame}

% SUMÁRIO

\begin{frame}{Sumário}
    \tableofcontents
\end{frame}

% INTRODUÇÃO

\section{Introdução}

\begin{frame}{Introdução}
    \begin{itemize}
        \item O trabalho apresenta um validador de placas veiculares brasileiras.
        
        \item O sistema reconhece dois padrões:
        \begin{itemize}
            \item Padrão antigo;
            \item Padrão Mercosul.
        \end{itemize}
        
        \item A solução utiliza Autômatos Finitos Determinísticos (AFD).
        
        \item A validação verifica somente o formato da placa, não sua existência em cadastros oficiais.
    \end{itemize}
\end{frame}

\begin{frame}{Objetivo do Projeto}
    \begin{block}{Objetivo principal}
        Desenvolver um validador capaz de identificar se uma placa
        veicular brasileira corresponde ao padrão antigo ou ao padrão
        Mercosul, utilizando conceitos da Teoria da Computação.
    \end{block}

    \vspace{0.4cm}

    \begin{itemize}
        \item Representar os formatos por linguagens formais.
        \item Construir autômatos para reconhecer os padrões.
        \item Demonstrar formalmente o funcionamento da solução.
        \item Analisar a minimalidade dos autômatos.
    \end{itemize}
\end{frame}

% FUNCIONAMENTO

\section{Como Funciona}

\begin{frame}{Funcionamento do Validador}
    O funcionamento consiste em receber uma placa e verificar
    se seus caracteres seguem algum dos padrões reconhecidos.

    \vspace{0.4cm}

    \begin{center}
    \begin{tikzpicture}[
        node distance=1.5cm,
        bloco/.style={
            rectangle,
            rounded corners,
            draw,
            align=center,
            minimum width=2.5cm,
            minimum height=0.9cm
        },
        >=Stealth
    ]

    \node[bloco] (entrada) {Entrada da placa};
    \node[bloco, right=of entrada] (afd) {Validação pelos AFDs};
    \node[bloco, right=of afd] (resultado) {Resultado};

    \draw[->] (entrada) -- (afd);
    \draw[->] (afd) -- (resultado);

    \end{tikzpicture}
    \end{center}

    \vspace{0.5cm}

    \begin{itemize}
        \item Se corresponder ao padrão antigo, será reconhecida.
        \item Se corresponder ao padrão Mercosul, será reconhecida.
        \item Caso contrário, será considerada inválida.
    \end{itemize}
\end{frame}

\begin{frame}{Limitações do Validador}
    \begin{alertblock}{Importante}
        O validador reconhece apenas a estrutura e o formato
        das placas.
    \end{alertblock}

    \begin{itemize}
        \item Não consulta bancos de dados oficiais.
        \item Não verifica se a placa está registrada.
        \item Não confirma a situação legal do veículo.
        \item Apenas identifica se a sequência de caracteres
        corresponde a um dos padrões definidos.
    \end{itemize}
\end{frame}

% ESPECIFICAÇÃO

\section{Especificação Formal}

\begin{frame}{Alfabeto}
    Para representar as placas, foram definidas as seguintes
    classes de símbolos:

    \vspace{0.3cm}

    \begin{itemize}
        \item $\Lt$: conjunto de letras de A até Z.
        \item $\Nm$: conjunto de algarismos de 0 até 9.
        \item $\Hf$: símbolo correspondente ao hífen.
    \end{itemize}

    \vspace{0.3cm}

    \begin{block}{Alfabeto}
        \[
            \Sigma = \Lt \cup \Nm \cup \Hf
        \]
    \end{block}
\end{frame}

\begin{frame}{Padrão Antigo}
    O padrão antigo de placas brasileiras possui:

    \begin{itemize}
        \item Três letras;
        \item Um hífen;
        \item Quatro números.
    \end{itemize}

    \begin{block}{Representação formal}
        \[
            \lang{Antigo} = \Lt^3 \Hf \Nm^4
        \]
    \end{block}

    \begin{exampleblock}{Exemplo}
        \centering
        \Large \texttt{ABC-1234}
    \end{exampleblock}
\end{frame}

\begin{frame}{Padrão Mercosul}
    O padrão Mercosul possui sete caracteres:

    \begin{itemize}
        \item Três letras iniciais;
        \item Um número na quarta posição;
        \item Uma letra na quinta posição;
        \item Dois números nas posições finais.
    \end{itemize}

    \begin{block}{Representação formal}
        \[
            \lang{Mercosul} =
            \Lt^3 \Nm \LAJ \Nm^2
        \]
    \end{block}

    \begin{exampleblock}{Exemplo}
        \centering
        \Large \texttt{ABC1D23}
    \end{exampleblock}
\end{frame}

\begin{frame}{Linguagem das Placas}
    A linguagem completa do projeto é formada pela união
    dos dois padrões.

    \begin{block}{Definição formal}
        \[
            \lang{Placas} =
            \lang{Antigo} \cup \lang{Mercosul}
        \]
    \end{block}

    \begin{itemize}
        \item Uma placa será aceita se pertencer a um dos dois padrões.
        \item Caso não pertença a nenhum deles, será rejeitada.
    \end{itemize}
\end{frame}

% MODELAGEM

\section{Modelagem dos Autômatos}

\begin{frame}{Autômato Finito Determinístico}
    Um AFD é definido formalmente pela tupla:

    \[
        M = (Q, \Sigma, \delta, q_0, F)
    \]

    Onde:

    \begin{itemize}
        \item $Q$: conjunto finito de estados.
        \item $\Sigma$: alfabeto de entrada.
        \item $\delta$: função de transição.
        \item $q_0$: estado inicial.
        \item $F$: conjunto de estados finais.
    \end{itemize}

    \vspace{0.3cm}

    O autômato percorre os estados de acordo com os
    caracteres recebidos e aceita a entrada quando
    termina em um estado final.
\end{frame}

\begin{frame}{Autômato do Padrão Antigo}
    \begin{center}
    \begin{tikzpicture}[
        shorten >=1pt,
        node distance=1.15cm,
        on grid,
        auto,
        ->,
        >=Stealth,
        every state/.style={minimum size=0.65cm}
    ]

    \node[state, initial] (q0) {$q_0$};
    \node[state, right=of q0] (q1) {$q_1$};
    \node[state, right=of q1] (q2) {$q_2$};
    \node[state, right=of q2] (q3) {$q_3$};
    \node[state, right=of q3] (q4) {$q_4$};
    \node[state, right=of q4] (q5) {$q_5$};
    \node[state, right=of q5] (q6) {$q_6$};
    \node[state, right=of q6] (q7) {$q_7$};
    \node[state, right=of q7, accepting] (q8) {$q_8$};

    \path
    (q0) edge node {$\Lt$} (q1)
    (q1) edge node {$\Lt$} (q2)
    (q2) edge node {$\Lt$} (q3)
    (q3) edge node {$\Hf$} (q4)
    (q4) edge node {$\Nm$} (q5)
    (q5) edge node {$\Nm$} (q6)
    (q6) edge node {$\Nm$} (q7)
    (q7) edge node {$\Nm$} (q8);

    \end{tikzpicture}
    \end{center}

    \vspace{0.3cm}

    \small
    O estado final é alcançado após a leitura das três letras,
    do hífen e dos quatro algarismos.
\end{frame}

\begin{frame}{Autômato do Padrão Mercosul}
    \begin{center}
    \begin{tikzpicture}[
        shorten >=1pt,
        node distance=1.15cm,
        on grid,
        auto,
        ->,
        >=Stealth,
        every state/.style={minimum size=0.65cm}
    ]

    \node[state, initial] (q0) {$q_0$};
    \node[state, right=of q0] (q1) {$q_1$};
    \node[state, right=of q1] (q2) {$q_2$};
    \node[state, right=of q2] (q3) {$q_3$};
    \node[state, right=of q3] (q4) {$q_4$};
    \node[state, right=of q4] (q5) {$q_5$};
    \node[state, right=of q5] (q6) {$q_6$};
    \node[state, right=of q6, accepting] (q7) {$q_7$};

    \path
    (q0) edge node {$\Lt$} (q1)
    (q1) edge node {$\Lt$} (q2)
    (q2) edge node {$\Lt$} (q3)
    (q3) edge node {$\Nm$} (q4)
    (q4) edge node {$\LAJ$} (q5)
    (q5) edge node {$\Nm$} (q6)
    (q6) edge node {$\Nm$} (q7);

    \end{tikzpicture}
    \end{center}

    \vspace{0.3cm}

    \small
    O autômato aceita a entrada quando os sete caracteres
    seguem a sequência definida pelo padrão Mercosul.
\end{frame}

% MINIMALIDADE

\section{Minimalidade}

\begin{frame}{Prova de Minimalidade}
    A minimalidade busca demonstrar que o autômato utiliza
    a menor quantidade possível de estados para reconhecer
    determinada linguagem.

    \vspace{0.3cm}

    \begin{itemize}
        \item Cada posição da placa exige uma etapa de leitura.
        \item Estados representam diferentes quantidades e tipos
        de caracteres já reconhecidos.
        \item Estados que representam situações distintas não
        podem ser simplesmente unidos.
    \end{itemize}

    \begin{block}{Objetivo}
        Demonstrar que não existe um AFD equivalente com
        uma quantidade menor de estados.
    \end{block}
\end{frame}

% CONCLUSÃO

\section{Conclusão}

\begin{frame}{Conclusão}
    \begin{itemize}
        \item O projeto utiliza conceitos de linguagens formais
        e autômatos finitos determinísticos.
        
        \item Foram considerados os padrões antigo e Mercosul.
        
        \item A modelagem permite reconhecer formalmente
        sequências válidas de caracteres.
        
        \item O estudo também aborda a definição formal
        e a minimalidade dos autômatos.
    \end{itemize}

    \vspace{0.4cm}

    \begin{block}{Resultado}
        A utilização de AFD permite implementar uma validação
        estruturada, determinística e eficiente dos formatos
        de placas veiculares brasileiras.
    \end{block}
\end{frame}

% ENCERRAMENTO

\begin{frame}{Obrigado!}
    \begin{center}
        \Huge Obrigado pela atenção!

        \vspace{0.5cm}

        \Large Perguntas?
    \end{center}
\end{frame}

\end{document}
